\section*{Author background and Section 3.1: Samplers}
Recall that a function $C:\Sigma^k\to\Sigma^n$ is an error-correcting code with distance $\delta$ if for every distinct $x,y\in\Sigma^k$, $\operatorname{dist}(C(x),C(y))\geq\delta$, where $\operatorname{dist}$ is the relative Hamming distance. The rate of the code $C$ is given by $\rho=k/n$. When there is more than one possible message length, we consider a code family, which is a family of functions in which each function is a code, and there is one code per message length $k$. A code family is asymptotically good if both the rate and distance of every code in the family are uniformly bounded below by constants $\rho>0$ and $\delta>0$, respectively.
Our distance amplification procedure makes use of samplers. These are bipartite graphs with a certain pseudo-random property. Let $G=(L,R,E)$ be a bipartite graph. We say $G$ is left-regular if all nodes in $L$ have the same degree.
\setcounter{nttheorem}{13}
\begin{ntdefinition}[Author Definition 14, {\cite{ref9}}]
Let $0<\varepsilon,\delta<1$. A bipartite graph $G=(L,R,E)$ is an $(\varepsilon,\delta)$-sampler if for every subset $T\subseteq R$, for all but $\delta$-fraction of vertices $v\in L$ it holds that
\[\left|\frac{|\Gamma(v)\cap T|}{|\Gamma(v)|}-\mu(T)\right|\leq\varepsilon.\]
\end{ntdefinition}
We will be working with ``unbalanced'' samplers. These are samplers with $|L|\gg|R|$. The state-of-the-art constructions of these samplers rely on their connection to randomness seeded extractors. We refer the interested reader to the excellent survey by Goldreich~\cite{ref17} for more information. When working with samplers, it is rather typical that the bipartite graph is left-regular, that is, the degree of all vertices in $L$ is the same. A small additional technical property we need is that the degree of every vertex in $R$ is close to the average right-degree. We make use of the following theorem which gives (non-explicit) samplers with near-optimal parameters having the above properties with respect to the degrees. We give a proof sketch for completeness.
\begin{nttheorem}[Author Theorem 15]
There exists a universal constant $c_{\mathrm{samp}}\geq1$ such that the following holds. For all integers $\ell,r$ and all $\varepsilon>0$, $1/2>\delta>0$ for which $\ell\geq r/(\delta\log(1/\delta))$, there exists a left-regular $(\varepsilon,\delta)$-sampler $G=([\ell],[r],E)$ with left-degree $d=c_{\mathrm{samp}}\cdot\log(1/\delta)/\varepsilon^2$. Moreover, provided that $\log r<1/(\delta\varepsilon^2)$, every right vertex has degree in $[D/2,2D]$ where $D=\ell d/r$ is the average right degree.
\end{nttheorem}
For the proof we need the following well-known lemma.
\begin{ntlemma}[Author Lemma 16]
For every integers $1\leq k\leq n$ with $k/n=\delta\leq1/2$ it holds that
\[\sum_{i=0}^k\binom ni\leq2^{H(\delta)n},\]
where $H(x)=-x\log(x)-(1-x)\log(1-x)$ is the binary entropy function.
\end{ntlemma}
\begin{proof}[Author proof sketch for Theorem 15]
The proof is via the probabilistic method, where for every left vertex we choose $d$ neighbours independently and uniformly at random, and independently across all left vertices (note that in the above we allow for parallel edges, but if that troubles the reader, that can be avoided as well in the regime of interest $d\ll r$ by arguing that the probability of a right neighbor to be selected more than once is small. In any case, our distance amplification procedure works just as well with parallel edges). Fix $T\subseteq[r]$. For $v\in[\ell]$ let $F_v$ be the indicator random variables that is $1$ if and only if $\bigl||\Gamma(v)\cap T|/d-\mu(T)\bigr|>\varepsilon$. By the Chernoff bound, $\Pr[F_v]\leq e^{-\Omega(\varepsilon^2d)}$. Fix $S\subseteq[\ell]$ with $|S|=\delta\ell$. The probability that for all vertices $v\in S$ it holds that $F_v=1$ is bounded above by $e^{-\Omega(\varepsilon^2d\cdot\delta\ell)}$. By taking the union bound over all $S$ and $T$, we get that except with probability
\begin{equation}2^r\binom\ell{\delta\ell}e^{-\Omega(\varepsilon^2d\delta\ell)}\leq2^{r+H(\delta)\ell-c\varepsilon^2d\delta\ell},\tag{3.1}\end{equation}
the sampled graph is an $(\varepsilon,\delta)$-sampler. Note that the last inequality follows by Lemma 16, where $c>0$ is some constant. By taking $c_{\mathrm{samp}}\geq5/c$, one can verify (using that $H(x)\leq2x\log(1/x)$ for all $x\leq1/2$) that the right hand side in Equation (3.1) is bounded by $1/4$.
As for the moreover part, again, by the Chernoff bound, the probability that there exists a right vertex which has degree outside $[D/2,2D]$ is bounded above by $re^{-\Omega(\ell d/r)}$, and this is bounded by $1/4$ by our choice of parameters and by taking $c_{\mathrm{samp}}$ large enough.
\end{proof}
\setcounter{nttheorem}{17}
\begin{nttheorem}[Author Theorem 18: The Gilbert--Varshamov bound]
Let $\Sigma$ be a set of size $|\Sigma|=q$. For every $n\in\mathbb N$, and $0\leq\delta\leq1-1/q$ there exists a code of block-length $n$ over $\Sigma$, with distance at least $\delta$ and rate $r\geq1-H_q(\delta)$. Furthermore, if $q$ is a prime power and $\Sigma=\mathbb F_q$, there exists a linear code over $\Sigma$ with rate $r\geq1-H_q(\delta)-g(n)$, where $g(n)=O(1/n)$.
\end{nttheorem}
